%!TEX root = ../template.tex
%%%%%%%%%%%%%%%%%%%%%%%%%%%%%%%%%%%%%%%%%%%%%%%%%%%%%%%%%%%%%%%%%%%
%% chapter1.tex
%% NOVA thesis document file
%%
%% Chapter with introduction
%%%%%%%%%%%%%%%%%%%%%%%%%%%%%%%%%%%%%%%%%%%%%%%%%%%%%%%%%%%%%%%%%%%

\typeout{NT FILE chapter1.tex}%

\chapter{Introduction}
\label{cha:introduction}

\prependtographicspath{{Chapters/Figures/Covers/}}

% epigraph configuration


\section{Motivation and Context}
\label{sec:if_you_use_this_template}

This first Chapter introduces the template and how it is organized. In Chapter~\ref{cha:users_manual} you can find some specific instructions on how to use this template.  Chapter~\ref{cha:a_short_latex_tutorial_with_examples} shows some examples and give some hints on how to write your text. Please read these next Chapters carefully.

\subsection{Motivation and Context}
\label{sub:Motivation_Context}

Did you learn how to drive by sitting by the wheel and throwing your car into the road?  Most probably you did take your time \emph{learning the rules} and \emph{practicing} first! Likewise, it is not wise to throw yourself at the task of writing a thesis/dissertation in \LaTeX\ without seriously considering the following \ntindex{recommendation}!

\begin{tcolorbox}[colback=green!8]
 Teste caso queira usar
\end{tcolorbox}

\subsection{Problem Definition}
\label{sub:prob_definition}
Problem definition


\subsection{Objectives}
\label{sub:objectives}
Thesis objectives

\subsection{Proposed Solution}
\label{sub:proposed_solution}
Proposed Solution